\textheight24cm
\topmargin-2.5cm
\oddsidemargin+0cm
\textwidth16.3cm 

\usepackage{amsmath}
\usepackage{amsfonts}
\usepackage{amssymb}
\usepackage{color}
\usepackage[latin1]{inputenc}
\usepackage{ifthen}
\usepackage{enumerate}
\usepackage{lastpage}
\usepackage{graphicx}
%\usepackage{subfigure}
\usepackage{subcaption}
\usepackage{bm}
\usepackage{tikz}
\usepackage{pgfplots}
\usepackage{float}
\usepackage{nicefrac}
\usepackage{epsfig}
\usepackage{etoolbox}
\usepackage{framed}
\usepackage{hyperref}
\usepackage{datetime}
\usepackage{url}
\usepackage{algorithm}
\usepackage{svg}

\usepackage[noend]{algpseudocode}
\algnewcommand{\parState}[1]{\State%
    \parbox[t]{\dimexpr\linewidth-\algmargin}{\strut\hangindent=\algorithmicindent \hangafter=1 #1\strut}}

\algrenewcommand\algorithmicindent{1.0em}%
\renewcommand\algorithmicdo{:}
\renewcommand\algorithmicthen{:}

%\newcommand{\linetofill}{\ \\\hphantom{\hspace{0mm}}\hrulefill\ \\}
\newcommand{\linetofill}{\ \\\hphantom{\hspace{0mm}}\dotfill\ \\}
\newcommand{\vlinetofill}[1]{\\\rule{#1 mm}{0.1mm}\ \\}
\newcommand{\cbox}{
  \setlength{\unitlength}{1pt}
  \begin{picture}(10,10)
    \put(0,0){\line(1,0){8}} \put(0,8){\line(1,0){8}}
    \put(0,0){\line(0,1){8}} \put(8,0){\line(0,1){8}}
  \end{picture}
}
\newcommand{\bigcbox}{
  \setlength{\unitlength}{3pt}
  \begin{picture}(10,10)(1,1)
    \put(0,0){\line(1,0){8}} \put(0,8){\line(1,0){8}}
    \put(0,0){\line(0,1){8}} \put(8,0){\line(0,1){8}}
  \end{picture}
}
% points and boxes on the right
\newcommand{\points}[1]{
\ \\[-5mm]
\hphantom{\ }\hfill\textbf{#1 pts \bigcbox \hspace{-6mm}}
\vspace{5mm}
}
% multiple choice checkboxes
%\newcommand{\boxt}{\hspace*{2em} $[~~]$ \textsf{True}}
%\newcommand{\boxf}{\hspace*{2em} $[~~]$ \textsf{False}}
\newcommand{\boxt}{\hspace*{2em} $\square$ \textsf{True}}
\newcommand{\boxf}{\hspace*{2em} $\square$ \textsf{False}}
\newcommand{\justif}{\textsf{Justification}:  \rule{0ex}{2em}\dotfill\\\rule{0ex}{2em}\dotfill}
\newcommand{\checkboxWithJustification}{\boxt \boxf\\ \justif\\}
\newcommand{\checkbox}{\boxt \boxf}
\usepackage{tikz}
\usepackage{xcolor}
\usepackage{amssymb}
\usepackage{amsmath}
\usepackage{amsthm}
\usepackage{pgfplots}
\pgfmathdeclarefunction{gauss}{2}{%
  \pgfmathparse{1/(#2*sqrt(2*pi))*exp(-((x-#1)^2)/(2*#2^2))}%
}
\renewcommand{\S}{{\cal S}}
\usepackage{pgfplots}
    \pgfplotsset{compat=newest}
    \pgfplotsset{small, every non boxed x axis/.append style={x axis line style=-},
every non boxed y axis/.append style={y axis line style=-}}

\usepackage{bbm}

\newcommand{\solution}[1]{{\color{ETHRed}{#1}}}
\NeedsTeXFormat{LaTeX2e}
\ProvidesPackage{llm}[2014/08/24 Example LaTeX package]

%%%%%%%%%%%%%%%%%%%%%%%%%%%%%%%%%%%%%%%%%%%%%%%%%%%%%%%
%%% COLORS %%%%%%%%%%%%%%%%%%%%%%%%%%%%%%%%%%%%%%%%%%%%
%%%%%%%%%%%%%%%%%%%%%%%%%%%%%%%%%%%%%%%%%%%%%%%%%%%%%%%
\usepackage{colortbl}

% ETH colors
\definecolor{ETHBlue}{RGB}{33,92,175}	% blue 
\definecolor{ETHGreen}{RGB}{98,115,19}		% green
\definecolor{ETHPurpleDark}{RGB}{140,10,89}	% purple
\definecolor{ETHPurple}{RGB}{163,7,116}	% purple
\definecolor{ETHGray}{RGB}{111,111,111}	% gray
\definecolor{ETHRed}{RGB}{183,53,45}	% red
\definecolor{ETHPetrol}{RGB}{0,120,148}	% green/blue 
\definecolor{ETHBronze}{RGB}{142,103,19}	% bronze

% ETH color aliases 
\colorlet{ETHdarkblue}{ETHBlue!80!black}
\colorlet{ETHdarkgreen}{ETHGreen!80!black}
\colorlet{ETHpink}{ETHPurple}
\colorlet{ETHgray}{ETHGray}
\colorlet{ETHred}{ETHRed}
\colorlet{ETHgreenblue}{ETHPetrol}
\colorlet{ETHbrown}{ETHBronze}

% Define the colors
\definecolor{TextBlack}{RGB}{51,51,51}
\definecolor{BackgroundWhite}{RGB}{255,255,255}

% Blues
\definecolor{AccentBlue}{RGB}{0,122,204}
\definecolor{LightBlue}{RGB}{173,216,230}
\definecolor{DarkBlue}{RGB}{0,51,102}

% Greens
\definecolor{AccentGreen}{RGB}{70,160,73}
\definecolor{LightGreen}{RGB}{144,238,144}
\definecolor{DarkGreen}{RGB}{0,100,0}

% Reds
\definecolor{AccentRed}{RGB}{255,0,0}
\definecolor{LightRed}{RGB}{255,99,71}
\definecolor{DarkRed}{RGB}{139,0,0}

% Oranges
\definecolor{AccentOrange}{RGB}{255,165,0}
\definecolor{LightOrange}{RGB}{255,204,153}
\definecolor{DarkOrange}{RGB}{255,140,0}

% Neutral Colors
\definecolor{NeutralLightGray}{RGB}{204,204,204}
\definecolor{NeutralMediumGray}{RGB}{102,102,102}

% Yellow for notes/highlights
\definecolor{NoteYellow}{RGB}{255,255,0}

% Diverse Colors
\definecolor{DiversePurple}{RGB}{128,0,128}
\definecolor{DiverseTeal}{RGB}{0,128,128}
\definecolor{DiverseOlive}{RGB}{128,128,0}
\definecolor{DiverseCyan}{RGB}{0,128,192}
\definecolor{DiverseMagenta}{RGB}{192,0,128}

%%%%%%%%%%%%%%%%%%%%%%%%%%%%%%%%%%%%%%%%%%%%%%%%%%%%%%%
%%% PACKAGES %%%%%%%%%%%%%%%%%%%%%%%%%%%%%%%%%%%%%%%%%%
%%%%%%%%%%%%%%%%%%%%%%%%%%%%%%%%%%%%%%%%%%%%%%%%%%%%%%%
\usepackage[english]{babel}

% \usepackage[margin=1in]{geometry}
\usepackage{geometry}

% \usepackage{mathpazo}

% math packages
\usepackage{amsthm}
\usepackage{amsmath}
\usepackage{amsfonts}
\usepackage{amssymb}
% scientific notation
\usepackage{siunitx}
\usepackage{etoolbox}


\newtoggle{publish-notes}
\iftoggle{publish-notes}{
    \colorlet{MacroColor}{black}
}{
    \colorlet{MacroColor}{ETHPetrol}
}
\colorlet{MACROCOLOR}{MacroColor}

\usepackage{natbib}
\newcommand*{\nameadjunct}{\relax}
\makeatletter
\renewcommand*{\NAT@nmfmt}[1]{\NAT@up #1\nameadjunct}
\makeatother
\newcommand*{\citeposs}[2][]{%
    \begingroup
    \renewcommand*{\nameadjunct}{'s}%
    \citet[#1]{#2}%
    \endgroup
}
\usepackage{xfrac}
\usepackage{mathtools}
\usepackage{bm}
\usepackage{bbm}
\usepackage{blkarray}
\usepackage{nicefrac}
\usepackage{xspace}
\usepackage[makeroom]{cancel}

\usepackage{datetime}

\usepackage{graphicx}
\usepackage{adjustbox}
\usepackage{relsize}
\usepackage{stmaryrd}

\usepackage{url}

\usepackage{booktabs}
\usepackage{multirow}

\usepackage{blindtext}

\usepackage{pdfpages}

\usepackage{microtype}
\usepackage{inconsolata}
\usepackage{scalerel}

\usepackage{thmtools}
\usepackage{thm-restate}

\usepackage{array}
\usepackage[inline]{enumitem}
\usepackage{caption}
\usepackage{subcaption}
\usepackage{listings}

\usepackage{mathabx}

\usepackage{imakeidx}
\makeindex[columns=2, options= -s indexstyle.ist]

\usepackage{titling}


\renewcommand{\qedsymbol}{$\blacksquare$}

% code rendering
%\usepackage[finalizecache,cachedir=.]{minted}
%\usepackage{textgreek}

% tikz packages
\usepackage{tikz}
\usetikzlibrary{automata, fit, positioning, matrix, arrows.meta, decorations.pathmorphing, shapes, shapes.multipart, chains,shadows.blur, shadows}
\usetikzlibrary{shadings}

% plotting packages
\usepackage{tikz-qtree}

\usepackage{tikz-dependency}

%tikz FSA settings
\tikzset{
transition/.style={-{Latex[length=1.8mm, width=1.4mm]}},
every loop/.append style=-{Latex[length=1.8mm, width=1.4mm]},
initial text=$ $,
node distance=2cm,
every state/.style={thick, draw=gray, fill=gray!15, circular drop shadow},
every initial by arrow/.style={text=ETHGreen,-{Latex[length=1.8mm, width=1.4mm]}},
initial text = $ $
}

%% Cleveref settings
\usepackage{cleveref}
\crefname{section}{\S}{\S\S}
\crefname{table}{Tab.}{Tabs.}
\crefname{figure}{Fig.}{Figs.}
\crefname{algorithm}{Algorithm}{Algorithms}
\crefname{equation}{Eq.}{Eqs.}
\crefname{fact}{Fact}{Facts}
\crefname{appendix}{Appendix}{Appendices}
\crefname{assumption}{Assumption}{Assumptions}
\crefname{lemma}{Lemma}{Lemmas}
\crefname{proposition}{Proposition}{Propositions}
\crefname{chapter}{Chapter}{Chapters}
\crefname{line}{line}{lines}
\crefname{principle}{Principle}{Principles}
\crefname{corollary}{Corollary}{Corollaries}
\crefname{Exercise}{Exercise}{Exercises}
\crefformat{section}{\S#2#1#3}

% Tokenization 
\DeclareCaptionType{myexamplef}[Example][List of Examples]
\crefname{myexamplef}{Example}{}
\DeclareCaptionType{codef}[Snippet][List of Snippets]
\crefname{codef}{Snippet}{}

% TODO notes
\setlength{\marginparwidth}{2cm}
\iftoggle{publish-notes}{
    \usepackage[disable]{todonotes}
}{
    \usepackage{todonotes}
}
\makeatletter
\newcommand*\iftodonotes{\if@todonotes@disabled\expandafter\@secondoftwo\else\expandafter\@firstoftwo\fi}  %
\makeatother
\newcommand{\noindentaftertodo}{\iftodonotes{\noindent}{}}
\newcommand{\fixme}[2][]{\todo[color=yellow,size=\scriptsize,fancyline,caption={},#1]{#2}}
% \renewcommand{\note}[4][]{\todo[author=#2,color=#3,size=\scriptsize,fancyline,caption={},#1]{#4}} % default note settings, used by macros below.
\newcommand{\note}[4][]{\todo[author=#2,color=#3,size=\scriptsize,fancyline,caption={},#1]{#4}} % default note settings, used by macros below.

\newcommand{\ryan}[2][]{\note[#1]{ryan}{violet!40}{#2}}
\newcommand{\clara}[2][]{\note[#1]{clara}{orange!40}{#2}}
\newcommand{\nik}[2][]{\note[#1]{nik}{yellow!40}{#2}}
\newcommand{\afra}[2][]{\note[#1]{afra}{coral}{#2}}
\newcommand{\anej}[2][]{\note[#1]{anej}{AccentGreen!80!LightGreen}{#2}}
\newcommand{\leo}[2][]{\note[#1]{leo}{cyan}{#2}}
\newcommand{\clemente}[2][]{\note[#1]{clemente}{green}{#2}}
\newcommand{\tianyu}[2][]{\note[#1]{tianyu}{red!40}{#2}}
\newcommand{\luca}[2][]{\note[#1]{luca}{magenta}{#2}}
\newcommand{\alexandra}[2][]{\note[#1]{alexandra}{teal!40}{#2}}
\newcommand{\franz}[2][]{\note[#1]{franz}{orange}{#2}}

\newcommand{\Ryan}[2][]{\ryan[inline,#1]{#2}\noindentaftertodo}
\newcommand{\Clara}[2][]{\clara[inline,#1]{#2}\noindentaftertodo}
\newcommand{\Anej}[2][]{\anej[inline,#1]{#2}\noindentaftertodo}
\renewcommand{\Leo}[2][]{\leo[inline,#1]{#2}\noindentaftertodo}
\newcommand{\Andreas}[2][]{\andreas[inline,#1]{#2}\noindentaftertodo}
\newcommand{\Eleanor}[2][]{\eleanor[inline,#1]{#2}\noindentaftertodo}
\newcommand{\Alexandra}[2][]{\alexandra[inline,#1]{#2}\noindentaftertodo}
\newcommand{\Tianyu}[2][]{\tianyu[inline,#1]{#2}\noindentaftertodo}
\newcommand{\Timv}[2][]{\timv[inline,#1]{#2}\noindentaftertodo}
\newcommand{\Clemente}[2][]{\clemente[inline,#1]{#2}\noindentaftertodo}
\newcommand{\Luca}[2][]{\luca[inline,#1]{#2}\noindentaftertodo}
\newcommand{\Franz}[2][]{\franz[inline,#1]{#2}\noindentaftertodo}

\newcommand{\response}[1]{\vspace{3pt}\hrule\vspace{3pt}\textbf{#1:}}    % insert \response{myname} within someone else's todonote

% algorithms stuff
\usepackage{algorithm}
\usepackage[noend]{algpseudocode}  % do not also use algorithmic
\algrenewcommand\algorithmicindent{1.0em}%
\renewcommand\algorithmicdo{:}
\renewcommand\algorithmicthen{:}
\newcommand{\rightcomment}[1]{{\color{gray} \(\triangleright\) {\footnotesize\textit{#1}}}}
\algrenewcommand{\algorithmiccomment}[1]{\hfill \rightcomment{#1}}  % redefines \Comment
\algnewcommand{\LineComment}[1]{\State\rightcomment{#1}}
\algnewcommand{\LinesComment}[1]{\State\rightcomment{\parbox[t]{.95\linewidth-\leftmargin-\widthof{\(\triangleright\) }}{#1}}}
\renewcommand\algorithmicdo{:}
\renewcommand\algorithmicthen{:}
\algrenewcommand\alglinenumber[1]{{\tiny\color{black!50}#1.}\hspace{-2pt}}

\newcommand{\algorithmicfunc}[1]{\textbf{def} {#1}:}
\algdef{SE}[FUNC]{Func}{EndFunc}[1]{\algorithmicfunc{#1}}{}
\makeatletter
\ifthenelse{\equal{\ALG@noend}{t}}%
{\algtext*{EndFunc}}
{}%
\makeatother
